\documentclass[11pt]{article}
\usepackage[margin=0.78in]{geometry}
\usepackage[T1]{fontenc}
\usepackage{lmodern}
\usepackage{microtype}
\usepackage{amsmath,amssymb,amsthm}
\usepackage{booktabs,tabularx,array}
\usepackage{enumitem}
\usepackage{fancyhdr}
\usepackage{xcolor}
\usepackage{hyperref}
\hypersetup{hidelinks}
\setlength{\parindent}{0pt}
\setlength{\parskip}{5.2pt}
\setlist[itemize]{leftmargin=1.35em,itemsep=2pt,topsep=3pt}
\setlist[enumerate]{leftmargin=1.55em,itemsep=2pt,topsep=3pt}
\newtheorem{proposition}{Proposition}[section]
\newtheorem{corollary}[proposition]{Corollary}
\newtheorem{definition}[proposition]{Definition}
\pagestyle{fancy}
\fancyhf{}
\lhead{\small LANGUAGE MECHANICS}
\rhead{\small MINIMUM FORMAL / MECHANICAL REFERENCE v0.2}
\cfoot{\thepage}
\renewcommand{\headrulewidth}{0.35pt}
\newcommand{\status}[1]{\textbf{Status.} #1}
\newcommand{\lawbox}[1]{\begin{center}\fcolorbox{black!60}{black!2}{\parbox{0.88\linewidth}{\centering #1}}\end{center}}
\begin{document}

\begin{center}
{\Large\bfseries 006 Minimum Formal / Mechanical Reference}\\[3pt]
{\large A Page-Zero Deck for Typed Mathematical Use}\\[10pt]
Anthony Vito Coppa\\[2pt]
\textit{Original reference constructions: 2026}\\
\textit{Canonical standalone edition v0.2: 29 August 2026}
\end{center}
\vspace{3pt}
\hrule
\vspace{8pt}

\status{\textsc{Declared} as a Language Mechanics reference instrument. Definitions are constitutive declarations of this deck; propositions and corollaries below are \textsc{Derived} from those declarations.}

\textbf{Abstract.} This deck fixes the minimum typing required to carry a formal reference through an operation without silently exchanging reader equality, carrier equality, history, lift, or refusal. It isolates a reference packet, typed return, receipt adequacy, recognition, the admission gate, and the distinction between \emph{minimal} and \emph{minimum}. A finite mechanical witness is included only after the operations it uses are declared: the seal carrier is a join-semilattice and the residue carrier is a monoid. The deck is locally complete and imports no physical, semantic, numerical, or cross-domain identification.

\section{Position}
This is a reference deck, not a replacement foundation for mathematics. It asks one operational question:

\lawbox{When a formal expression carries a referent through an operation, what must be declared so that the result can be compared with its reference without type loss?}

The deck begins from ordinary formal objects, operations, equality relations, and maps. It adds only explicit typing of reference, return, readable difference, lift, receipt, recognition, minimum, and refusal.

\section{The minimum reference packet}
Fix a jurisdiction
\[
D=(\Sigma,R,\rho),
\]
where \(\Sigma\) is the candidate class, \(R\) is the permitted re-posing relation, and \(\rho:\Sigma\to Y\) is the declared reader.

\begin{definition}[Reference packet]
A formal reference packet is
\[
\mathfrak R=(D,x_0,T,\rho),
\]
where \(x_0\in\Sigma\) is the addressed referent and \(T:\operatorname{Dom}(T)\subseteq\Sigma\to\Sigma\) is a declared partial operation.
\end{definition}

The elementary transport is
\[
x_0\xrightarrow{\,T\,}T(x_0),
\]
and the first return question is
\[
\rho(Tx_0)\stackrel{?}{=}\rho(x_0).
\]
No statement about identity, memory, meaning, or mechanism follows until its own office is declared.

\section{Holding-through and typed equality}
\begin{definition}[Reference survives operation]
The reference survives \(T\) at \(x_0\) under reader \(\rho\) when
\[
\rho(Tx_0)=\rho(x_0).
\]
If the full carrier is the reader, \(\rho=\operatorname{id}_\Sigma\), this reduces to exact carrier return \(Tx_0=x_0\).
\end{definition}

\begin{proposition}[Typed equality]\label{prop:typed-equality}
From
\[
\rho(Tx)=\rho(x)
\]
one may conclude only that \(Tx\) and \(x\) lie in the same fiber of \(\rho\). Carrier equality \(Tx=x\) follows only when the declared hypotheses make that fiber a singleton for the compared pair---for example, when \(\rho\) is injective there.
\end{proposition}
\begin{proof}
Equality of the reads says \(Tx\in\rho^{-1}(\rho(x))\). Without injectivity, a fiber may contain distinct points, so equality in \(Y\) does not force equality in \(\Sigma\). If \(\rho\) is injective on the relevant domain, the implication follows by definition of injectivity.
\end{proof}

\status{Proposition~\ref{prop:typed-equality} is \textsc{Derived}.}

\section{Minimum office register}
The following table is an entrance index. Any stronger use of a term requires the stronger structure to be declared where it is used.

\small
\begin{tabularx}{\textwidth}{@{}>{\bfseries}p{0.88in}X X@{}}
\toprule
Office & Exact minimum use & Boundary\\
\midrule
Jurisdiction & \(D=(\Sigma,R,\rho)\): where a claim is posed & does not supply content\\
Reference & addressed comparison baseline & does not explain why selected\\
Condition & entry requirement & does not guarantee continuation\\
Constraint & relation preserved by admitted continuation & distinct from condition\\
\(\mathsf{can}\) & lawful continuation in the declared jurisdiction & gate verdict, not truth\\
\(\varnothing\) & refused / undefined continuation here & not scalar zero\\
Operation & typed map or partial map & no semantics supplied by the map alone\\
Difference & distinguishable departure under shared reference & requires a comparison frame\\
Return & equality at a declared reader & not automatically carrier reset\\
Lift & finer carrier/address above a quotient or read & requires an explicit map\\
Receipt & recoverable evidence of occurrence at promised resolution & not the state itself\\
Recognition & warranted present comparison under a declared reader & not unrestricted identity\\
Minimum & least admitted construction under a declared order & stronger than mere minimality\\
Fit & witnessed compatibility preserving the required relation & resemblance is insufficient\\
Function & operation lawfully carried through admitted Fit & does not create Fit\\
\bottomrule
\end{tabularx}
\normalsize

\section{Exact return, reader return, and lift}
Two predicates must remain separate:
\[
\text{exact carrier return:}\qquad Tx=x,
\]
\[
\text{reader return with retained carrier difference:}\qquad
\rho(Tx)=\rho(x),\quad Tx\neq x.
\]
The second form is the minimum opening for a lift: the declared read closes while a finer carrier distinction remains.

\begin{definition}[Lift datum]
Given a quotient or reader \(p:\widetilde X\to X\), a lift datum for a returned action consists of a lifted action \(\widetilde T\) and a witness
\[
p(\widetilde T\widetilde x)=p(\widetilde x),
\]
with, whenever nontriviality is claimed,
\[
\widetilde T\widetilde x\neq\widetilde x.
\]
\end{definition}
The quotient return and the lifted difference are separate reads of one declared construction.

\section{Receipt and recognition}
A return may erase information from its displayed state. Receipt therefore occupies a separate office.

\begin{definition}[Receipt adequacy]\label{def:receipt}
Let \(H\) be a history class and let \(\sim_\delta\) be the equivalence relation induced by the promised history resolution \(\delta\), with quotient map
\[
\pi_\delta:H\to H/{\sim_\delta}.
\]
Let \(\Sigma_\delta\) be a receipt space and
\[
\sigma_\delta:H\to\Sigma_\delta
\]
a receipt map. The receipt is adequate at resolution \(\delta\) when the promised quotient remains recoverable. A sufficient exact condition is a decoder
\[
\operatorname{Dec}_\delta:\operatorname{im}(\sigma_\delta)\to H/{\sim_\delta}
\]
satisfying
\[
\boxed{\operatorname{Dec}_\delta\circ\sigma_\delta=\pi_\delta.}
\]
\end{definition}
The two sides now have the same domain \(H\); the equation states exactly that the receipt retains every history distinction promised by the quotient.

\begin{definition}[Recognition]
Recognition at reader \(\rho\) is a warranted present comparison
\[
\operatorname{Rec}_\rho(x;x_0)=1
\]
only when the declared comparison is defined, the stated equality or bounded relation holds, and the evidence required by the claim is present.
\end{definition}
Recognition returns only what its reader and receipt can carry.

\section{The gate and refusal}
For a candidate move \(T\) at \(x\), let \(C(x)\) be its entry condition and let \(\mathcal Q(x,Tx)\) be the continuation constraint. Define
\[
\operatorname{Adm}(T,x)
\iff
C(x)\wedge (T\text{ is typed at }x)\wedge\mathcal Q(x,Tx),
\]
and
\[
\operatorname{Gate}(T,x)=
\begin{cases}
\mathsf{can},&\operatorname{Adm}(T,x),\\
\varnothing,&\text{otherwise}.
\end{cases}
\]
The namespace is locked:
\[
\boxed{\varnothing\neq0.}
\]
A refused continuation is not a scalar, zero vector, empty history, false proposition, or absent physical object.

\section{Minimal and minimum}
The word \emph{minimum} carries a stronger burden than \emph{minimal}.

\begin{definition}[Minimal and minimum admitted constructions]
Let \((\mathcal A,\preceq)\) be a declared partially ordered class of admitted constructions, with strict part \(\prec\).
A construction \(m\in\mathcal A\) is \emph{minimal} when
\[
\nexists c\in\mathcal A\text{ such that }c\prec m.
\]
It is a \emph{minimum} when
\[
\boxed{m\preceq c\qquad\text{for every }c\in\mathcal A.}
\]
\end{definition}

\begin{proposition}[Minimum implies minimal]
Every minimum is minimal. A minimal element need not be a minimum.
\end{proposition}
\begin{proof}
If \(m\preceq c\) for every \(c\in\mathcal A\), no \(c\prec m\) can exist. Conversely, two incomparable minimal elements supply a class with minimal elements but no minimum.
\end{proof}

Therefore every formal use of \emph{minimum} must declare at least:
\begin{itemize}
\item the admitted candidate class;
\item the admission conditions;
\item the comparison/reduction order;
\item the relations that must be preserved;
\item the proof that the selected construction lies below every admitted competitor when a true minimum is claimed.
\end{itemize}
Without those declarations, ``minimum'' is not yet a formal result.

\section{Mechanical seal/residue witness}
A finite witness can separate cumulative reference from carried non-closing difference, but its operations must be typed before use.

Let \((\mathsf S,\vee)\) be a finite join-semilattice with its induced order \(\leq\). Let \((A,\oplus_A,0_A)\) be a monoid. Define
\[
\Gamma=\mathsf S\times A.
\]
For a declared turn \(k\), choose \(e_k\in\mathsf S\) and \(\epsilon_k\in A\), and define
\[
U_k(\lambda,\eta)=\bigl(\lambda\vee e_k,\;\eta\oplus_A\epsilon_k\bigr).
\]

\begin{corollary}[Well-defined monotone seal update]
Each \(U_k:\Gamma\to\Gamma\) is well defined, and its seal coordinate is monotone:
\[
\lambda\leq\lambda\vee e_k.
\]
\end{corollary}
\begin{proof}
The join exists by the join-semilattice law, and \(\oplus_A\) is closed on \(A\) by the monoid law. The defining property of a join makes \(\lambda\vee e_k\) an upper bound of \(\lambda\).
\end{proof}

The coordinates have different jobs:
\[
\boxed{\text{seal records admitted cumulative structure}}
\qquad
\boxed{\text{residue carries non-closing difference}}.
\]
This witness imposes no historical twelve-coordinate semantic map, spectral/color read, distinguished numerical constant, calendar, or physical realization.

\section{Page-zero rule and boundary}
\lawbox{\textbf{A formal term must survive the operation under which it is reused.}}

Expanded:
\[
\text{reference}\to\text{typed operation}\to\text{return test}\to\text{receipt}\to\text{recognition}\to\text{reuse}.
\]
This is a discipline of formal use. It does not claim that mathematics, physics, language, or cognition are one mechanism. It does not import certification from a neighboring paper or from repeated vocabulary. Any stronger construction must state its own carrier, operations, reader, preserved relation, witness, and falsifier.

\section*{Source record}
The following sources are provenance only; no theorem above depends on an unquoted result from them.
\begin{itemize}
\item A. V. Coppa, \emph{Rule \& Reference: Page Zero}, 10 June 2026.
\item A. V. Coppa, \emph{Classical Reference: A Return Typing of Classical Mathematics}, working paper, 2026.
\item A. V. Coppa, \emph{Minimum Mechanical Reference: Seal / Residue Form}, working paper, 2026.
\end{itemize}

\end{document}
